\documentclass[10pt,a4paper]{amsart}
\usepackage{amsmath,amssymb,amsthm,graphicx}
\usepackage{amsmath}
\usepackage{amssymb}
\usepackage{hyperref}
% Minimal class-compatibility preamble for a reader excerpt of
%   Juan Esteban Rodriguez Camargo, "An integral model of the perfectoid modular curve",
%   Journal de l'Ecole polytechnique -- Mathematiques 8 (2021), DOI 10.5802/jep.170.
%
% The block below is the author's own macro and theorem declaration block of the pinned
% source file, copied mechanically and unchanged (source lines 6-91), followed by the
% abbreviations that the publisher class cedram supplies and that the retained author text
% uses (same meanings as in that class). No mathematical macro, statement, symbol or
% assumption is introduced, changed or removed by this file.
\usepackage{xspace}

\usepackage[scr=rsfs,cal=euler]{mathalfa}
\multlinegap0pt

\newenvironment{enumeratei}
{\bgroup\def\theenumi{\roman{enumi}}\def\theenumii{\arabic{enumii}}\begin{enumerate}}
{\end{enumerate}\egroup}

\newcommand\bw[1]{\mathbin{\hbox{$\mathsf{\Lambda}$}^{\!#1}}}
\newcommand{\quand}{\quad\text{and}\quad}
\newcommand{\sbullet}{{\scriptscriptstyle\bullet}}

\newcommand\mto{\mathchoice{\longmapsto}{\mapsto}{\mapsto}{\mapsto}}
\newcommand\hto{\mathrel{\lhook\joinrel\to}}
\newcommand\To[1]{\mathchoice{\xrightarrow{\textstyle\kern4pt#1\kern3pt}}{\stackrel{#1}{\to}}{}{}}
\let\ra\rightarrow

\let\ts\textstyle
\newcommand{\lcr}{[\![}
\newcommand{\rcr}{]\!]}
\newcommand{\lpr}{(\!(}
\newcommand{\rpr}{)\!)}

\newcommand{\rae}{\mathrm{ae}}
\newcommand{\an}{\mathrm{an}}
\newcommand{\rb}{\mathrm{b}}
\newcommand{\rc}{\mathrm{c}}
\newcommand{\qc}{\mathrm{qc}}
\newcommand{\cusp}{\mathrm{cusp}}
\newcommand{\cusps}{\mathrm{cusps}}
\newcommand{\cyc}{\mathrm{cyc}}
\newcommand{\et}{\mathrm{\acute{e}t}}
\newcommand{\fTd}{\mathrm{fTd}}
\newcommand{\pr}{\mathrm{pr}}
\newcommand{\proet}{\mathrm{pro}\textup{-}\et}
\newcommand{\sm}{\mathrm{sm}}
\newcommand{\univ}{\mathrm{univ}}

\newcommand{\wotimes}{\mathbin{\widehat\otimes}}

\usepackage[all,cmtip]{xy}\let\labelstyle\textstyle
\usepackage{tikz-cd}

\numberwithin{equation}{section}

\newtheorem{prop}{Proposition}[section]
\newtheorem{lem}[prop]{Lemma}
\newtheorem{theo}[prop]{Theorem}
\newtheorem{cor}[prop]{Corollary}

\theoremstyle{definition}
\newtheorem*{claim}{Claim}
\newtheorem{exam}[prop]{Example}
\newtheorem{definition}[prop]{Definition}
\newtheorem{remark}[prop]{Remark}

\newcommand{\bb}[1]{\mathbb{#1}}
\newcommand{\n}[1]{\mathcal{#1}}
\newcommand{\f}[1]{\mathfrak{#1}}
\newcommand{\s}[1]{\mathscr{#1}}

\DeclareMathOperator{\GL}{GL}
\DeclareMathOperator{\Tr}{Tr}
\DeclareMathOperator{\id}{id}
\DeclareMathOperator{\SL}{SL}
\DeclareMathOperator{\R}{R}
\DeclareMathOperator{\D}{D}
\DeclareMathOperator{\N}{N}

\DeclareMathOperator{\Hom}{Hom}
\DeclareMathOperator{\dotimes}{\otimes^{\mathrm{L}}}
\DeclareMathOperator{\coker}{coker}

\DeclareMathOperator{\Gal}{Gal}
\DeclareMathOperator{\Spf}{Spf}
\DeclareMathOperator{\Spec}{Spec}
\DeclareMathOperator{\Ell}{Ell}
\DeclareMathOperator{\Spa}{Spa}
\DeclareMathOperator{\Cusps}{Cusps}
\DeclareMathOperator{\Tate}{Tate}
\DeclareMathOperator{\sHom}{\s{H}\!\mathit{om}}
\DeclareMathOperator{\HomSurj}{HomSurj}
\DeclareMathOperator{\Art}{Art}
\DeclareMathOperator{\Intperf}{\f{I}nt\textup{-}\f{P}erf}
\DeclareMathOperator{\Perf}{\s{P}erf}
\newcommand{\mc}[1]{X(Np^{#1})}
\newcommand{\fmc}[1]{\mathfrak{X}(Np^{#1})}

% Publisher-class abbreviations (cedram) used by the retained author text.
\providecommand{\loccit}{\emph{loc. cit.}}
\providecommand{\bibliofont}{\small}
\providecommand{\backmatter}{}
\providecommand{\skpt}{\smallskip}
\providecommand{\smaller}{\small}
\newcommand{\ie}{i.e.,\xspace}
\newcommand{\eg}{e.g.,\xspace}
\newcommand{\etc}{etc.\xspace}
\newcommand{\cf}{cf.\xspace}
\newcommand{\resp}{resp.\xspace}
\title{The upper shriek of the structure sheaf along an lci morphism is a shifted invertible sheaf}
\author{Juan Esteban Rodríguez Camargo}
\date{doi:10.5802/jep.170; CC BY 4.0}
\begin{document}
\maketitle
\noindent\textit{Editorial note.} All delivered mathematical text is the author's own, copied from the pinned publisher source sources/algebra-fl03-an-integral-model-of-the-perfectoi/pinned-original-source.tex; no statement or proof step was rewritten, reordered or removed.\par\medskip
\noindent\textit{Editorial note.} Preamble provenance: the retained preamble is the source's own retained-preamble.tex, which carries the author's macro and theorem declaration block verbatim (its own header records which lines of the pinned file it copies); the wrapper class supplies the baseline packages that block assumes. Nothing was added, no macro was invented, and the excerpt serves only as preamble, never as a span coordinate.\par\medskip
\noindent\textit{Editorial note.} Metadata provenance: source\_id and source\_version were taken from sources/algebra-fl03-an-integral-model-of-the-perfectoi/license-observation.json; the title from SOURCE.json; the author list from the source metadata record; the source URL from the first non-empty of source\_url, source\_urls[0] and the evidence's fixed\_version\_url/license\_url. The licence evidence file used by this card is sources/algebra-fl03-an-integral-model-of-the-perfectoi/license-observation.json, recorded in license.evidence\_path so that the binding target is explicit.\par\medskip
\noindent\textit{Editorial note.} External references are navigation pointers only; every deliberately kept step of the author's argument stays inside the counted spans. The displayed numbers of those pointers are the author's own numbering under the author's \textbackslash{}newtheorem declarations (section.index as declared), computed mechanically by the lane's scanner.\par\medskip
\noindent\textit{Editorial note.} The counted proof contains no reference to any labelled item and no bibliography citation, so the card needs neither a navigation stub nor a bibliography excerpt.\par\medskip
\noindent\textit{Editorial note.} The remark that follows the lemma refers back to it; that reference lies outside the counted ranges and is therefore not part of the delivered material.\par\medskip
\noindent\textit{Allowed standard prerequisites (editorial).}\begin{itemize}
\item The local factorization of an lci morphism into a smooth morphism followed by a regular closed immersion. Used inside the counted proof at lines 694-695. Reduces the statement to a smooth morphism g and a regular closed immersion iota, for which the upper shriek is computed separately.
\item The composition formula for the upper shriek and the identification of g\textasciicircum{}! of the structure sheaf with shifted m-differentials for a smooth morphism. Used inside the counted proof at lines 696-700. Produces the displayed chain of equalities for f\textasciicircum{}! of the structure sheaf.
\item The Koszul complex of a regular sequence is a flat resolution, and its dual resolves the quotient shifted by the codimension. Used inside the counted proof at lines 701-707. Replaces the derived hom by a tensor computation and identifies the result with the restriction of the invertible sheaf of m-differentials.
\end{itemize}\medskip
\begin{lem}
\label{fOy}
Let $f: X\to Y$ be an lci morphism of relative dimension $n$ between locally noetherian schemes of finite Krull dimension. Then $f^!\s{O}_Y = \omega^{\circ}_{X/Y}[n]$ with $\omega^{\circ}_{X/Y}$ an invertible $\s{O}_X$-module.
\end{lem}

\begin{proof}
Working locally on $Y$ and $X$, we may assume that $f$ factors as $X \To{\iota} S \To{g} Y$, where $g$ is a smooth morphism of relative dimension $m$, and $\iota$ is a regular closed immersion of codimension $m-n$ defined by an ideal $\s{I}=(f_1,\ldots, f_{m-n})$. Let $\omega^{\circ}_{S/Y}=\bw{m}\Omega^{1}_{S/Y}$ be the sheaf of $m$-differentials of $S$ over $Y$, then
\begin{align*}
f^!\s{O}_Y & = \iota^! g^! \s{O}_Y \\
		 & = \iota^{-1} \R \sHom_{\s{O}_S}(\iota_* \s{O}_X, \omega_{S/Y}^\circ[m]) \\
		 & = \iota^{-1} \R \sHom_{\s{O}_S}(\s{O}_S/\s{I}, \omega_{S/Y}^\circ)[m].
\end{align*}
Let $K(\underline{f})$ be the Koszul complex of the regular sequence $\underline{f}=(f_1,\ldots, f_{m-n})$. Then $K(\underline{f})$ is a flat resolution of $\s{O}_S/\s{I}$, its dual $K(\underline{f})^\vee=\sHom_{\s{O}_S}(K(\underline{f}), \s{O}_S)$ is a flat resolution of \mbox{\((\s{O}_S/\s{I})[-(m-n)]\)}. Therefore
\begin{align*}
f^!\s{O}_Y & \simeq \iota^{-1} \sHom_{\s{O}_S}(K(\underline{f}), \omega^\circ_{S/Y})[m] \\
		 & = \iota^{-1} K(\underline{f})^\vee \otimes \omega^\circ_{S/Y} [m] \\
		 & \simeq \iota^{-1}(\s{O}_{S}/\s{I}\otimes ^L \omega^\circ_{S/Y}[n]) \\
		 & = \iota^{-1}((\omega^\circ_{S/Y}/\s{I})[n]) = \iota^* \omega^{\circ}_{S/Y}[n],
\end{align*}
which is an invertible sheaf of $\s{O}_X$-modules as required.
\end{proof}

\begin{thebibliography}{99}
\bibitem{hartshorne1966residues} R. Hartshorne - Residues and duality , Lect. Notes in Math. , vol. 20 , Springer-Verlag, Berlin-New York, 1966
\bibitem{mumford1974abelian} D. Mumford - Abelian varieties , TIFR Studies in Mathematics , vol. 5 , Tata Institute of Fundamental Research, Bombay, 2008
\end{thebibliography}
\end{document}
